\begin{center}
{\Large\bfseries 5. Rationale Funktionenkörper.\srcfn{1)}{Vortrag, gehalten auf der Naturforscherversammlung Wien 1913.}}\par
\vspace{1.2em}
Von \textsc{Emmy Noether} in Erlangen.\par
\medskip
J. Ber. d. DMV 22 (1913), S. 316--319
\end{center}

\vspace{2em}

Die folgenden Fragestellungen gehen ursprünglich auf Gespräche mit E. Fischer zurück. Einige der Fragen sind übrigens für den speziellen Fall einer Unbestimmten --- und zwar unter allgemeineren Voraussetzungen über den Koeffizientenbereich --- schon von E. Steinitz aufgeworfen und erledigt.\srcfn{2)}{E. Steinitz: Algebraische Theorie der Körper (besonders \S\ 24). Crelles Journal 137. 1910.}

Unter einem \glqq{}rationalen Funktionenkörper\grqq{} verstehe ich einen Körper, dessen Elemente rationale Funktionen von $n$ Unbestimmten sind, mit Koeffizienten aus einem vorgegebenen Zahlkörper, der insbesondere auch alle komplexen Zahlen umfassen kann. Beispiele solcher Körper bilden der Körper der symmetrischen Funktionen von $n$ Größen, oder allgemeiner die Gesamtheit der rationalen Funktionen von $n$ Unbestimmten, die die Permutationen einer gewissen Gruppe gestatten (Lagrangesche Gattungsbereiche); ferner der Invariantenkörper.

Zwischen den beiden erstgenannten Körpern und dem letzteren besteht ein charakteristischer Unterschied: die ersteren enthalten $n$ algebraisch unabhängige Funktionen, die symmetrischen Elementarfunktionen; während die Anzahl der algebraisch unabhängigen Invarianten stets kleiner ist als die Anzahl der Unbestimmten. Es ist deshalb wichtig, daß man allgemein solche Körper zweiter Art eineindeutig den Körpern erster Art zuordnen kann, indem man einige der Unbestimmten durch Zahlen ersetzt. Ich will mich also im folgenden auf Körper erster Art, d. h. auf Körper mit $n$ algebraisch unabhängigen Funktionen beschränken; die Resultate gelten vermöge der Zuordnung dann allgemein.

Die Fragestellungen gruppieren sich um drei Basisbegriffe: Rationalbasis, Minimalbasis und Integritätsbasis.

1. Unter einer \glqq{}Rationalbasis\grqq{} verstehe ich eine endliche Anzahl von Funktionen des Körpers derart, daß jede Funktion des Körpers sich als rationale Verbindung dieser endlichen Anzahl darstellen läßt, mit Koeffizienten aus dem vorgegebenen Koeffizientenbereich. Man zeigt durch einfache Überlegungen --- die im Fall einer Unbestimmten sich auch bei E. Steinitz a. a. O. finden --- die Existenz der Rationalbasis für jeden rationalen Funktionenkörper. Beispielsweise bilden nach dem Lagrangeschen Satz die symmetrischen Elementarfunktionen und eine Funktion, die zur Gruppe gehört, eine Rationalbasis für die früher erwähnten Lagrangeschen Gattungsbereiche.

2. Jede Rationalbasis muß (für Körper erster Art) mindestens $n$ Funktionen enthalten; enthält sie genau $n$ Funktionen, die dann algebraisch unabhängig sein müssen, so bezeichne ich sie als \glqq{}Minimalbasis\grqq{}. Die Frage nach der Existenz der Minimalbasis läßt sich teilweise beantworten durch Sätze von Lüroth, Castelnuovo und Enriques.\srcfn{1)}{J. Lüroth: Beweis eines Satzes über rationale Kurven. Math. Ann. 9. 1875. --- G. Castelnuovo: Sulla razionalità delle involuzioni piane. Math. Ann. 44. 1893. --- F. Enriques: Sopra una involuzione non razionale dello spazio. Rend. Acc. Linc. vol. 21. 21. Jan. 1912.} Danach existiert für Körper einer Unbestimmten stets eine Minimalbasis: die bis auf lineargebrochene Transformation eindeutig bestimmte Lürothsche Funktion des Körpers (vgl. Steinitz \S\ 24). Für Körper von zwei Unbestimmten existiert eine Minimalbasis stets dann und im allgemeinen nur dann, wenn man algebraische Erweiterung des Koeffizientenbereichs zuläßt, während für drei und mehr Unbestimmte eine Minimalbasis im allgemeinen überhaupt nicht existiert. Die speziellen Körper mit Minimalbasis zu charakterisieren, ist noch nicht versucht worden.

Ich habe nun die Frage nach der Minimalbasis bei den Lagrangeschen Gattungsbereichen weiter verfolgt. Hier gewinnt sie die folgende Bedeutung: \glqq{}Wenn der zur Gruppe $G$ gehörende Lagrangesche Gattungsbereich eine Minimalbasis besitzt, so kann man die Gesamtheit der Gleichungen mit vorgeschriebener Gruppe $G$ rational durch Parameterdarstellung konstruieren; und zwar für jeden Zahlkörper, der den Koeffizientenbereich der Minimalbasis enthält --- also insbesondere für jeden beliebigen Zahlkörper, wenn dieser Koeffizientenbereich der Körper der rationalen Zahlen ist.\grqq{}

Das einfachste Beispiel bietet die symmetrische Gruppe; hier bilden die symmetrischen Elementarfunktionen eine Minimalbasis mit rationalen Zahlkoeffizienten. Daraus ergibt sich als Parameterdarstellung einfach die Gleichung mit unbestimmten Koeffizienten. Wie man von dieser zu beliebig vielen affektlosen Gleichungen für jeden Zahlkörper übergeht, zeigt der Hilbertsche Irreduzibilitätssatz.

Nun ergibt sich direkt aus den Sätzen von Lüroth und Castelnuovo die Existenz der Minimalbasis für alle bei den Gleichungen 3. und 4. Grades auftretenden Gruppen; wobei sich weiter zeigt, daß diese Minimalbasis rationale Zahlkoeffizienten besitzt. Man kann also für jeden beliebigen Zahlkörper die Gesamtheit der Gleichungen 3. und 4. Grades mit vorgeschriebener Gruppe rational konstruieren. Für die zyklischen und Diedergruppen existiert eine Minimalbasis mit dem Körper der $n$ten Einheitswurzeln als Koeffizientenbereich; das gleiche gilt für einige weitere metazyklische Gruppen. Die Frage, ob die Minimalbasis überhaupt für gewisse Gruppen charakteristisch ist, muß unentschieden bleiben.

3. Um zu dem letzten Basisbegriff zu gelangen, betrachte ich die Gesamtheit der in dem Körper enthaltenen Polynome. Unter einer \glqq{}Integritätsbasis\grqq{} verstehe ich eine endliche Anzahl dieser Polynome derart, daß jedes Polynom des Körpers sich als ganze rationale Verbindung dieser endlichen Anzahl darstellen läßt, mit Koeffizienten aus dem vorgegebenen Koeffizientenbereich. Die Frage nach der Existenz der Integritätsbasis --- die z. B. als Spezialfall die Endlichkeit des Invariantensystems in sich schließt --- ist in etwas anderer Fassung von Hilbert in seinen \glqq{}Mathematischen Problemen\grqq{}\srcfn{1)}{D. Hilbert: Mathematische Probleme. Vortrag Paris 1900. Problem 14. Göttinger Nachr. 1900.} als Problem der relativ ganzen Funktionen aufgeworfen.

Ich kann einstweilen nur eine Klasse von Körpern angeben, für die eine Integritätsbasis existiert. Enthält nämlich der Körper ein System von $n$ Polynomen derart, daß die homogene Resultante der Glieder höchster Dimension dieser Polynome von Null verschieden ist, so existiert eine Integritätsbasis; diese ist gegeben durch die Koeffizienten aller $u$ der im Körper irreduzibeln Gleichung für die Linearform:
\[
  u_0=u_1x_1+u_2x_2+\cdots+u_nx_n.\srcfnmark{2)}
\]
\srcfntext{2)}{Diese irreduzible Gleichung findet sich im Falle einer Unbestimmten bei E. Steinitz zum Beweis des Lürothschen Satzes.}
Ist insbesondere der Wert der Resultante gleich einer Einheit des Koeffizientenbereichs, so ist auch die Ganzzahligkeit der Darstellung gesichert. Ein Beispiel zu dieser Klasse von Körpern bilden die Lagrangeschen Gattungsbereiche, da die Resultante der symmetrischen Elementarfunktionen den Wert $\pm 1$ hat. Ein weiteres Beispiel (ohne Ganzzahligkeit) ist gegeben durch alle Körper, die nur ein algebraisch unabhängiges Polynom enthalten; die Integritätsbasis ist hier identisch mit der Lürothschen Funktion des kleinsten, alle Polynome enthaltenden Teilkörpers. So sind bekanntlich alle ganzen rationalen, projektiven Invarianten einer quadratischen Form von $n$ Variabeln ganze Funktionen der Diskriminante.

Die angegebene Bedingung ist nur hinreichend, keineswegs notwendig für die Existenz der Integritätsbasis. Man kann leicht Körper konstruieren, bei denen die Resultante von je $n$ Polynomen verschwindet, und die dennoch eine durch die Koeffizienten jener irreduzibeln Gleichung gegebene Integritätsbasis besitzen. Es entsteht die Frage, ob etwa auch im allgemeinsten Fall die Koeffizienten jener Gleichung die Integritätsbasis liefern.
